% Reviewed translation: PDF pages 201--206 / printed pages 187--192.
\MathBookSourcePage{201}
为强调与连续情形的对应关系, 定义
\[
a_h(g_h,l_h)=\nu(\operatorname{grad}\mathbf u_h(g_h),
\operatorname{grad}\mathbf u_h(l_h))
\qquad\forall g_h,l_h\in G_h.
\]
显然, 问题 \eqref{eq:primitive-gp-discrete-zero-problems} 有唯一解; 对给定的
$g_h$, 问题 \eqref{eq:primitive-gp-discrete-harmonic-problems} 也有唯一解.
此外, 容易验证问题 \eqref{eq:primitive-gp-discrete-boundary-problem} 同样有唯一解
$\rho_h$. 事实上, 若 $\mathbf u_h(\rho_h)=\mathbf0$, 则
\[
(\operatorname{grad}p_h(\rho_h),\mathbf v_h)=0
\qquad\forall\mathbf v_h\in X_h.
\]
但 Corollary~\ref{cor:primitive-hood-taylor-global-inf-sup} 建立的 inf-sup 条件蕴含
\[
p_h(\rho_h)=\frac1{\operatorname{meas}(\Omega)}
\int_\Omega p_h(\rho_h)\,dx,
\]
即 $p_h(\rho_h)$ 在 $\Omega$ 上为常数. 又因为
$\int_\Gamma p_h(\rho_h)\,ds=0$, 所以 $p_h(\rho_h)=0$, 特别地
$\rho_h=0$. 因而得到下列结果.

\setcounter{statement}{2}
\renewcommand{\theHstatement}{lemma.\thesection.\arabic{statement}}
\begin{Lemma}\label{lem:primitive-gp-discrete-wellposedness}
设右端项 $\mathbf f\in L^2(\Omega)^2$, 其中 $\Omega$ 是有界平面多边形,
并假设三角剖分 $\mathcal T_h$ 满足
\eqref{eq:primitive-hood-taylor-macroelement-partition}. 则 Glowinski--Pironneau
格式 \eqref{eq:primitive-gp-discrete-zero-problems}--
\eqref{eq:primitive-gp-discrete-solution} 确定唯一的
$(\mathbf u_h,p_h)$, 其中 $\mathbf u_h\in X_h$, $p_h\in Q_h$, 且
$\int_\Gamma p_h\,ds=0$.

此外, $(\mathbf u_h,p_h)$ 满足
\begin{equation}\label{eq:primitive-gp-discrete-combined-equation}
\nu(\operatorname{grad}\mathbf u_h,\operatorname{grad}\mathbf v_h)
+(\operatorname{grad}p_h,\mathbf v_h-\operatorname{grad}q_h)
=(\mathbf f,\mathbf v_h-\operatorname{grad}q_h)
\quad\forall(\mathbf v_h,q_h)\in X_h\times\Phi_h.
\end{equation}
\end{Lemma}

注意, \eqref{eq:primitive-gp-discrete-combined-equation} 相当于两个相互独立的方程:
一个关于 $\mathbf u_h$, 另一个关于 $p_h$. 它们分别由合并
\eqref{eq:primitive-gp-discrete-zero-problems} 与
\eqref{eq:primitive-gp-discrete-harmonic-problems} 的最后一个方程和第一个方程得到.

需要指出的是, 尽管有限元空间与 Hood--Taylor 方法中的空间相同, 且两种情形
都满足取 $q_h=0$ 的 \eqref{eq:primitive-gp-discrete-combined-equation}, 上述
$(\mathbf u_h,p_h)$ 一般不是 Hood--Taylor 算法的解, 因为它不满足离散无散度约束
\[
(\mathbf u_h,\operatorname{grad}q_h)=0
\qquad\forall q_h\in Q_h.
\]
事实上, \eqref{eq:primitive-gp-discrete-harmonic-problems} 给出
\[
\begin{aligned}
(\operatorname{grad}p_h(g_h),\mathbf u_h)
&=-\nu(\operatorname{grad}\mathbf u_h(g_h),
\operatorname{grad}\mathbf u_h)\\
&=-\nu(\operatorname{grad}\mathbf u_h(g_h),
\operatorname{grad}(\mathbf u_h^0+\mathbf u_h(\rho_h)))\\
&=-(\mathbf u_h^0,\operatorname{grad}p_h(g_h))
-\nu(\operatorname{grad}\mathbf u_h(g_h),
\operatorname{grad}\mathbf u_h^0),
\end{aligned}
\]
其中最后一步使用了 \eqref{eq:primitive-gp-discrete-boundary-problem}.
再次应用 \eqref{eq:primitive-gp-discrete-harmonic-problems} 可得
\[
(\operatorname{grad}p_h(g_h),\mathbf u_h)=0
\qquad\forall g_h\in G_h.
\]
不同于连续情形, 该等式未必推广到 $Q_h$ 中所有 $q_h$. 从这个意义上说,
Glowinski--Pironneau 算法放松了无散度约束. 实际上, 为满足这一约束,
必须给 $\mathbf u_h$ 加上一个修正项.

\MathBookSourcePage{202}
具体地, 在 $\Phi_h$ 中定义 $\lambda_h$:
\begin{equation}\label{eq:primitive-gp-discrete-velocity-correction}
(\operatorname{grad}\lambda_h,\operatorname{grad}q_h)
=(\mathbf u_h,\operatorname{grad}q_h)
\qquad\forall q_h\in\Phi_h.
\end{equation}
则 $\mathbf u_h-\operatorname{grad}\lambda_h$ 满足
\[
(\mathbf u_h-\operatorname{grad}\lambda_h,\operatorname{grad}q_h)=0
\qquad\forall q_h\in Q_h.
\]
事实上, 取满足 $\int_\Gamma q_h\,ds=0$ 的 $q_h\in Q_h$, 并令
$g_h\in G_h$ 表示 $q_h$ 的边界值, 即 $q_h|_\Gamma=g_h|_\Gamma$.
于是, 关于 $|\cdot|_{1,\Omega}$, $q_h$ 有正交分解
\[
q_h=p_h(g_h)+q_h^0,
\qquad q_h^0\in\Phi_h,
\]
其中 $p_h(g_h)$ 是 \eqref{eq:primitive-gp-discrete-harmonic-problems} 的解. 因而
\[
\begin{aligned}
(\mathbf u_h,\operatorname{grad}q_h)
&=(\mathbf u_h,\operatorname{grad}q_h^0)\\
&=(\operatorname{grad}\lambda_h,\operatorname{grad}q_h^0)
&&\text{由 \eqref{eq:primitive-gp-discrete-velocity-correction}}\\
&=(\operatorname{grad}\lambda_h,\operatorname{grad}q_h)
&&\text{由 \eqref{eq:primitive-gp-discrete-harmonic-problems}}.
\end{aligned}
\]
下面将看到 $\lambda_h$ 确实很小, 因而 $\mathbf u_h$ 几乎“无散度”. 此外,
虽然这里只为理论目的引入 $\lambda_h$, 它在实际中高效求解
\eqref{eq:primitive-gp-discrete-boundary-problem} 时也很有用.

\setcounter{statement}{4}
\renewcommand{\theHstatement}{remark.\thesection.\arabic{statement}}
\begin{Remark}\label{rem:primitive-gp-direct-three-field-formulation}
三元组 $(\mathbf u_h,p_h,\lambda_h)$ 也可直接定义为
$X_h\times(Q_h/\mathbb R)\times\Phi_h$ 中下列问题的唯一解:
\[
\left\{
\begin{aligned}
\nu(\operatorname{grad}\mathbf u_h,\operatorname{grad}\mathbf v_h)
+(\operatorname{grad}p_h,\mathbf v_h-\operatorname{grad}q_h)
&=(\mathbf f,\mathbf v_h-\operatorname{grad}q_h)\\
&\quad\forall(\mathbf v_h,q_h)\in X_h\times\Phi_h,\\
(\mathbf u_h-\operatorname{grad}\lambda_h,\operatorname{grad}q_h)&=0
\quad\forall q_h\in Q_h.
\end{aligned}
\right.
\]
但公式 \eqref{eq:primitive-gp-discrete-zero-problems}--
\eqref{eq:primitive-gp-discrete-solution} 的优点是, 它们把解表示为 Laplace 算子
一系列解耦 Dirichlet 问题的解.
\end{Remark}

误差分析与相应问题
\eqref{eq:primitive-continuous-pressure-discrete-problem} 的误差分析非常相似.
特别地, 它的 inf-sup 条件恰为
\eqref{eq:primitive-continuous-pressure-inf-sup}, 因此 Theorem
\ref{thm:primitive-hood-taylor-local-inf-sup} 及其 Corollary 均适用.

\setcounter{statement}{5}
\renewcommand{\theHstatement}{theorem.\thesection.\arabic{statement}}
\begin{Theorem}\label{thm:primitive-gp-error}
设 $\Omega$ 是有界凸平面多边形, 且 Stokes 问题
\eqref{eq:primitive-continuous-pressure-stokes-system} 的右端项
$\mathbf f\in L^2(\Omega)^2$. 若该问题的解 $(\mathbf u,p)$ 满足
\[
\mathbf u\in H^{k+1}(\Omega)^2,
\qquad p\in H^k(\Omega)\cap L_0^2(\Omega),
\qquad k=1,2,
\]
且三角剖分 $\mathcal T_h$ 一致正则并满足
\eqref{eq:primitive-hood-taylor-macroelement-partition}, 则有误差估计
\begin{equation}\label{eq:primitive-gp-primary-errors}
|\mathbf u-\mathbf u_h|_{1,\Omega}
+|\lambda_h|_{1,\Omega}
+\|p-\bar p_h\|_{0,\Omega}
\leq C_1h^k\{\|\mathbf u\|_{k+1,\Omega}+|p|_{k,\Omega}\},
\end{equation}
\MathBookSourcePage{203}
\begin{equation}\label{eq:primitive-gp-pressure-h1-error}
|p-p_h|_{1,\Omega}
\leq C_2h^{k-1}\{\|\mathbf u\|_{k+1,\Omega}+|p|_{k,\Omega}\},
\end{equation}
\begin{equation}\label{eq:primitive-gp-corrected-velocity-l2-error}
\|\mathbf u-\mathbf u_h+\operatorname{grad}\lambda_h\|_{0,\Omega}
\leq C_3h^{k+1}\{\|\mathbf u\|_{k+1,\Omega}+|p|_{k,\Omega}\},
\end{equation}
其中 $(\mathbf u_h,p_h)$ 是
\eqref{eq:primitive-gp-discrete-zero-problems}--
\eqref{eq:primitive-gp-discrete-solution} 的解, $\lambda_h$ 由
\eqref{eq:primitive-gp-discrete-velocity-correction} 给出, $\bar p_h$ 是
$p_h$ 在 $L_0^2(\Omega)$ 中的代表元.
\end{Theorem}

\begin{Proof}
有
\begin{equation}\label{eq:primitive-gp-error-equations}
\left\{
\begin{aligned}
\nu(\operatorname{grad}(\mathbf u-\mathbf u_h),\operatorname{grad}\mathbf v_h)
+(\operatorname{grad}(p-p_h),\mathbf v_h-\operatorname{grad}q_h)&=0\\
&\quad\forall(\mathbf v_h,q_h)\in X_h\times\Phi_h,\\
(\mathbf u_h-\operatorname{grad}\lambda_h,
\operatorname{grad}q_h)&=0
\quad\forall q_h\in Q_h.
\end{aligned}
\right.
\end{equation}
把 $(\mathbf v_h,q_h)$ 限制在空间
\[
B_h=\{(\mathbf v_h,q_h)\in X_h\times\Phi_h;
(\mathbf v_h-\operatorname{grad}q_h,\operatorname{grad}\mu_h)=0
\ \forall\mu_h\in Q_h\}.
\]
注意 $(\mathbf u_h,\lambda_h)\in B_h$. 于是
\eqref{eq:primitive-gp-error-equations} 写成
\[
\nu(\operatorname{grad}(\mathbf u-\mathbf u_h),\operatorname{grad}\mathbf v_h)
+(\operatorname{grad}(p-\mu_h),
\mathbf v_h-\operatorname{grad}q_h)=0
\]
对所有 $(\mathbf v_h,q_h)\in B_h$ 和所有 $\mu_h\in Q_h$ 成立.
为消去 $\operatorname{grad}q_h$, 取
$\mu_h=P_hp$, 其中 $P_hp$ 是由
\eqref{eq:appendix-neumann-ritz-projection} 定义的 $p$ 在 $M_h$ 上的
$H^1$ 投影. 因而
\[
\nu(\operatorname{grad}(\mathbf u-\mathbf u_h),\operatorname{grad}\mathbf v_h)
=(\operatorname{div}\mathbf v_h,p-P_hp)
\qquad\forall(\mathbf v_h,q_h)\in B_h.
\]
由于 $(\mathbf u_h,\lambda_h)\in B_h$, 该式立即给出
\[
|\mathbf u-\mathbf u_h|_{1,\Omega}
\leq2|\mathbf u-\mathbf w_h|_{1,\Omega}
+\frac{\sqrt2}{\nu}\|p-P_hp\|_{0,\Omega}
\qquad\forall(\mathbf w_h,q_h)\in B_h.
\]
这里显然可以取 $(\mathbf w_h,0)$, 其中 $\mathbf w_h\in V_h$.
由于 $(X_h,M_h)$ 满足 inf-sup 条件
\eqref{eq:primitive-continuous-pressure-inf-sup}, 可应用
\eqref{eq:primitive-discrete-affine-best-approximation}:
\[
\inf_{\mathbf w_h\in V_h}|\mathbf u-\mathbf w_h|_{1,\Omega}
\leq(1+\sqrt2/\beta^*)
\inf_{\mathbf v_h\in X_h}|\mathbf u-\mathbf v_h|_{1,\Omega}.
\]
所以
\[
|\mathbf u-\mathbf u_h|_{1,\Omega}
\leq2(1+\sqrt2/\beta^*)
\inf_{\mathbf v_h\in X_h}|\mathbf u-\mathbf v_h|_{1,\Omega}
+\frac{\sqrt2}{\nu}\|p-P_hp\|_{0,\Omega}.
\]
\eqref{eq:primitive-gp-primary-errors} 中的速度界由
\eqref{eq:primitive-hood-taylor-interpolation-error} 和
\eqref{eq:appendix-ritz-projection-lp-error} 得到. 注意, 只有
\eqref{eq:appendix-ritz-projection-lp-error} 需要 $\Omega$ 的凸性和
$\mathcal T_h$ 的一致正则性.

$\lambda_h$ 的界由上式和 $\operatorname{div}\mathbf u=0$ 得到:
\begin{equation}\label{eq:primitive-gp-correction-identity}
(\operatorname{grad}\lambda_h,\operatorname{grad}q_h)
=-(\mathbf u-\mathbf u_h,\operatorname{grad}q_h)
\qquad\forall q_h\in Q_h.
\end{equation}
因此
\[
|\lambda_h|_{1,\Omega}\leq\|\mathbf u-\mathbf u_h\|_{0,\Omega}.
\]

关于压力, 在 \eqref{eq:primitive-gp-error-equations} 中取 $q_h=0$,
即可回到标准情形. 因而由 inf-sup 条件和
\eqref{eq:primitive-pressure-error},
\MathBookSourcePage{204}
\[
\|p-\bar p_h\|_{0,\Omega}
\leq(1+\sqrt2/\beta^*)\inf_{q_h\in M_h}\|p-q_h\|_{0,\Omega}
+(\nu/\beta^*)|\mathbf u-\mathbf u_h|_{1,\Omega}.
\]
这给出 \eqref{eq:primitive-gp-primary-errors}; 再用熟知的论证得到
\eqref{eq:primitive-gp-pressure-h1-error}.

最后建立速度的 $L^2$ 估计. 证明是 Theorem
\ref{thm:primitive-duality-error} 证明的简单变体. 因为 $\Omega$ 为凸集,
存在唯一的
$(\boldsymbol\phi,\mu)\in[V\cap H^2(\Omega)^2]
\times[H^1(\Omega)\cap L_0^2(\Omega)]$ 使
\begin{equation}\label{eq:primitive-gp-dual-problem}
\begin{aligned}
(\mathbf g,\mathbf u-\mathbf u_h)
&=\nu(\operatorname{grad}\boldsymbol\phi,
\operatorname{grad}(\mathbf u-\mathbf u_h))
+(\operatorname{grad}\mu,\mathbf u-\mathbf u_h),\\
\|\boldsymbol\phi\|_{2,\Omega}+|\mu|_{1,\Omega}
&\leq C\|\mathbf g\|_{0,\Omega}.
\end{aligned}
\end{equation}
把该等式与 \eqref{eq:primitive-gp-error-equations} 结合, 得
\[
\begin{aligned}
(\mathbf g,\mathbf u-\mathbf u_h)
={}&\nu(\operatorname{grad}(\boldsymbol\phi-\boldsymbol\phi_h),
\operatorname{grad}(\mathbf u-\mathbf u_h))\\
&+(p-p_h,\operatorname{div}(\boldsymbol\phi_h-\boldsymbol\phi))
-(\mu-q_h,\operatorname{div}(\mathbf u-\mathbf u_h))\\
&+(\operatorname{grad}\lambda_h,
\operatorname{grad}(\mu-q_h))
-(\mathbf g,\operatorname{grad}\lambda_h)
\end{aligned}
\]
对所有 $\boldsymbol\phi_h\in X_h$, $q_h\in Q_h$ 成立
(这里使用 $\Delta\mu=\operatorname{div}\mathbf g$ in $\Omega$).
因此取 $q_h=P_h\mu$, 得
\[
\begin{aligned}
(\mathbf g,\mathbf u-\mathbf u_h+\operatorname{grad}\lambda_h)
={}&\nu(\operatorname{grad}(\boldsymbol\phi-\boldsymbol\phi_h),
\operatorname{grad}(\mathbf u-\mathbf u_h))\\
&+(p-p_h,\operatorname{div}(\boldsymbol\phi_h-\boldsymbol\phi))
-(\mu-P_h\mu,\operatorname{div}(\mathbf u-\mathbf u_h))
\end{aligned}
\]
对所有 $\boldsymbol\phi_h\in X_h$ 成立. 根据
\eqref{eq:primitive-gp-primary-errors} 和
\eqref{eq:primitive-gp-dual-problem}, 这给出
\eqref{eq:primitive-gp-corrected-velocity-l2-error}.
\end{Proof}

\setcounter{statement}{5}
\renewcommand{\theHstatement}{remark.\thesection.\arabic{statement}}
\begin{Remark}\label{rem:primitive-gp-no-uncorrected-l2-bound}
似乎不可能只对 $\mathbf u-\mathbf u_h$ 建立类似
\eqref{eq:primitive-gp-corrected-velocity-l2-error} 的 $L^2$ 界. 相反,
\eqref{eq:primitive-gp-correction-identity} 蕴含
\[
\|\mathbf u-\mathbf u_h+\operatorname{grad}\lambda_h\|_{0,\Omega}^2
=\|\mathbf u-\mathbf u_h\|_{0,\Omega}^2-|\lambda_h|_{1,\Omega}^2,
\]
因而不能从 \eqref{eq:primitive-gp-corrected-velocity-l2-error} 中单独分离出
$\mathbf u-\mathbf u_h$ 或 $\lambda_h$. 这并不令人惊讶, 因为 $\lambda_h$
充当速度 $\mathbf u_h$ 的修正.
\end{Remark}

当 $X_h$ 由 \eqref{eq:primitive-hood-taylor-refined-velocity-space} 定义
(而 $M_h$ 不变)时, 可以应用相同的分析. 由于
\eqref{eq:primitive-hood-taylor-refined-interpolation-error}, Theorem
\ref{thm:primitive-gp-error} 的陈述只对 $k=1$ 成立.

\subsection{Glowinski--Pironneau 格式的实现}
\label{subsec:primitive-gp-implementation}

计算 \eqref{eq:primitive-gp-discrete-zero-problems},
\eqref{eq:primitive-gp-discrete-harmonic-problems} 和
\eqref{eq:primitive-gp-discrete-boundary-problem} 的解并不完全直接, 因为测试函数
$l_h$ 没有显式出现在 \eqref{eq:primitive-gp-discrete-boundary-problem} 左端.
更容易的做法是通过计算
\eqref{eq:primitive-gp-discrete-velocity-correction} 定义的辅助函数 $\lambda_h$
来分解计算; 该函数补偿 $\mathbf u_h$ 不满足离散无散度约束这一事实.

\MathBookSourcePage{205}
具体地, 仿照 \eqref{eq:primitive-gp-discrete-zero-problems} 和
\eqref{eq:primitive-gp-discrete-harmonic-problems}, 定义 $\lambda_h(g_h)$ 与
$\lambda_h^0$:
\begin{equation}\label{eq:primitive-gp-implementation-corrections}
\left\{
\begin{aligned}
\lambda_h^0\in\Phi_h,\quad
(\operatorname{grad}\lambda_h^0,\operatorname{grad}q_h)
&=(\mathbf u_h^0,\operatorname{grad}q_h)
&&\forall q_h\in\Phi_h,\\
\lambda_h(g_h)\in\Phi_h,\quad
(\operatorname{grad}\lambda_h(g_h),\operatorname{grad}q_h)
&=(\mathbf u_h(g_h),\operatorname{grad}q_h)
&&\forall q_h\in\Phi_h.
\end{aligned}
\right.
\end{equation}
于是 $\lambda_h=\lambda_h^0+\lambda_h(\rho_h)$, 并且双线性形式
$a_h(\cdot,\cdot)$ 有另一种表达式.

\setcounter{statement}{3}
\renewcommand{\theHstatement}{lemma.\thesection.\arabic{statement}}
\begin{Lemma}\label{lem:primitive-gp-implementation-bilinear-form}
有
\begin{equation}\label{eq:primitive-gp-implementation-bilinear-form}
a_h(g_h,l_h)
=(\operatorname{grad}\lambda_h(g_h)-\mathbf u_h(g_h),
\operatorname{grad}l_h)
\qquad\forall g_h,l_h\in G_h.
\end{equation}
类似地, \eqref{eq:primitive-gp-discrete-boundary-problem} 的右端写成
\begin{equation}\label{eq:primitive-gp-implementation-right-hand-side}
(\mathbf u_h^0,\operatorname{grad}p_h(l_h))
=-(\operatorname{grad}\lambda_h^0-\mathbf u_h^0,
\operatorname{grad}l_h)
\qquad\forall l_h\in G_h.
\end{equation}
\end{Lemma}

\begin{Proof}
根据定义和 \eqref{eq:primitive-gp-discrete-harmonic-problems},
\[
\begin{aligned}
a_h(g_h,l_h)
&=\nu(\operatorname{grad}\mathbf u_h(g_h),
\operatorname{grad}\mathbf u_h(l_h))\\
&=-(\operatorname{grad}p_h(l_h),\mathbf u_h(g_h))\\
&=-(\operatorname{grad}(p_h(l_h)-l_h),\mathbf u_h(g_h))
-(\operatorname{grad}l_h,\mathbf u_h(g_h))\\
&=-(\operatorname{grad}\lambda_h(g_h),
\operatorname{grad}(p_h(l_h)-l_h))
-(\operatorname{grad}l_h,\mathbf u_h(g_h)),
\end{aligned}
\]
其中最后一步使用 \eqref{eq:primitive-gp-implementation-corrections}.
再由 \eqref{eq:primitive-gp-discrete-harmonic-problems},
\[
a_h(g_h,l_h)
=(\operatorname{grad}\lambda_h(g_h)-\mathbf u_h(g_h),
\operatorname{grad}l_h),
\]
从而证明 \eqref{eq:primitive-gp-implementation-bilinear-form}.
\eqref{eq:primitive-gp-implementation-right-hand-side} 的证明类似.
\end{Proof}

因此问题 \eqref{eq:primitive-gp-discrete-boundary-problem} 具有更易处理的形式:
求 $\rho_h\in G_h$ 使
\begin{equation}\label{eq:primitive-gp-implementation-boundary-problem}
(\operatorname{grad}\lambda_h(\rho_h)-\mathbf u_h(\rho_h),
\operatorname{grad}l_h)
=(\operatorname{grad}\lambda_h^0-\mathbf u_h^0,
\operatorname{grad}l_h)
\qquad\forall l_h\in G_h.
\end{equation}
由前一 Lemma 和双线性形式 $a_h(\cdot,\cdot)$ 的定义可知,
\eqref{eq:primitive-gp-implementation-boundary-problem} 左端是
$G_h\times G_h$ 上的双线性、对称、正定形式. 下面简述实际中如何求解
\eqref{eq:primitive-gp-implementation-boundary-problem}; 更多细节可见
Glowinski et al. 的 Chapter 13.

设 $\mathcal T_h\cap\Gamma$ 的节点编号为 $1$ 到 $N_h$, 并令
$\{\mu_i\}_{1\leq i\leq N_h}$ 为 $Q_h$ 的一组基函数, 定义为
\[
\mu_i(a_i)=1\quad\text{对 }\mathcal T_h\cap\Gamma\text{ 的每个节点 }a_i,
\qquad
\mu_i(b_k)=0\quad\text{对 }\mathcal T_h\text{ 的所有其他节点}.
\]
暂时略去条件 $\int_\Gamma g_h\,ds=0$, 则 $G_h$ 中函数可表示为
\MathBookSourcePage{206}
\[
g_h=\sum_{i=1}^{N_h}g_i\mu_i,
\qquad g_i=g_h(a_i).
\]
用这一记号, \eqref{eq:primitive-gp-implementation-boundary-problem} 等价于
\[
\sum_{j=1}^{N_h}\rho_j a_h(\mu_j,\mu_i)
=-(\operatorname{grad}\lambda_h^0-\mathbf u_h^0,
\operatorname{grad}\mu_i)
\qquad1\leq i\leq N_h.
\]
换言之, 需要求解线性方程组
\begin{equation}\label{eq:primitive-gp-implementation-linear-system}
A_h\boldsymbol\rho=\mathbf b,
\end{equation}
其中
\[
(A_h)_{i,j}=a_h(\mu_j,\mu_i)
=(\operatorname{grad}\lambda_h(\mu_j)-\mathbf u_h(\mu_j),
\operatorname{grad}\mu_i),
\]
\[
b_i=-(\operatorname{grad}\lambda_h^0-\mathbf u_h^0,
\operatorname{grad}\mu_i).
\]
为计算 $\mathbf b$, 需要求解
\eqref{eq:primitive-gp-discrete-zero-problems} 的三个 Dirichlet 问题以得到
$\mathbf u_h^0$, 再求解
\eqref{eq:primitive-gp-implementation-corrections} 的第一个 Dirichlet 问题以得到
$\lambda_h^0$, 共四个 Dirichlet 问题. 为计算矩阵 $A_h$ 的第 $j$ 列,
必须在 $g_h=\mu_j$ 时求解
\eqref{eq:primitive-gp-discrete-harmonic-problems} 的三个 Dirichlet 问题以得到
$\mathbf u_h(\mu_j)$, 再求解
\eqref{eq:primitive-gp-implementation-corrections} 的第二个问题以得到
$\lambda_h(\mu_j)$, 同样共四个 Dirichlet 问题.

由上述讨论, 矩阵 $A_h$ 对称且半正定, 并以 $0$ 为单重特征值. 此外,
若 $\mathcal T_h$ 的节点适当编号, 可以证明
$\operatorname{Ker}(A_h)$ 是常向量空间, 且主块
\[
\widetilde A_h=(a_h(\mu_j,\mu_i))_{1\leq i,j\leq N_h-1}
\]
正定. 因此令 $\widetilde\rho_{N_h}=0$, 可用 Cholesky 分解求解
\eqref{eq:primitive-gp-implementation-linear-system}, 得到 $\rho$ 的一个代表元
$\widetilde\rho$ 的前 $N_h-1$ 个分量. 随后, 满足
$\int_\Gamma\rho_h\,ds=0$ 的
\eqref{eq:primitive-gp-implementation-boundary-problem} 的解为
\[
\rho_h=\widetilde\rho_h-\frac1{\operatorname{meas}(\Gamma)}
\int_\Gamma\widetilde\rho_h\,ds.
\]
一旦 $\rho_h$ 已知, 即可通过求解
\eqref{eq:primitive-gp-discrete-harmonic-problems} 的三个 Dirichlet 问题计算压力
$p_h(\rho_h)$ 和速度 $\mathbf u_h(\rho_h)$.

问题 \eqref{eq:primitive-gp-implementation-boundary-problem} 也可用共轭梯度算法求解
(参见 Glowinski et al. 的上述文献).
